% GENERATED FILE. Edit canonical lessons, then run npm run generate:books.
% canonical-source-hash: fnv1a64:f4d3c1dd600bf98c
% canonical-lessons: GU-C310-tukdi, GU-C310-jivlen, GU-C310-sahay, GU-C310-mujhavu, GU-C310-vikheravu

\chapter{The Search Team}
\label{ch:gu-the-search-team}
\hlchaptermodality{310}

\begin{chapteropening}
\textbf{By the end of this chapter:} \emph{I can say how a team searched the forest.}

Complete the last lesson of chapter 310: I can say how a team searched the forest.
\end{chapteropening}

\section[\texorpdfstring{\gu{ટુકડી} (ṭukḍī)}{ટુકડી (ṭukḍī)}]{\gu{ટુકડી} (ṭukḍī) — a squad; a small party (of police or soldiers)}

\label{lesson:GU-C310-tukdi}



\begin{quote}
Before the new one: say the Gujarati for to applaud, then the Gujarati for pride.
\end{quote}

\subsection*{You'll want to know: \gu{ટુકડી}}
\textbf{\gu{ટુકડી}} — \emph{ṭukḍī} — \textquotedblleft{}a squad; a small party (of police or soldiers)\textquotedblright{}.

\begin{quote}
\gu{પોલીસની} \gu{એક} \gu{ટુકડી} \gu{જંગલમાં} \gu{ગઈ}.
\end{quote}

\begin{quote}
\emph{polīsnī ek ṭukḍī jangalmā̃ gaī.}
\end{quote}

\begin{quote}
A police squad went into the forest.
\end{quote}

A feminine noun.

\begin{cousinweb}
\textbf{\gu{ટુકડી}} (\emph{ṭukḍī}) is built on \emph{ṭukḍo}, a piece, which CDIAL lists under \emph{*ṭukk}, \textquotedblleft{}small\textquotedblright{}, a form it reconstructs (CDIAL 5466): a small piece of a larger body of people.
\end{cousinweb}

\begin{grammarlens}[title={What you've built}]
One more word for telling what happened.
\end{grammarlens}

\subsection*{Your turn}
\begin{itemize}
\raggedright
  \item \emph{Say it:} \emph{ṭukḍī}
  \item \emph{Say it:} \emph{ṭukḍī}, once more
  \item \emph{Say it:} \emph{polīsnī ek ṭukḍī jangalmā̃ gaī.} — A police squad went into the forest.
  \item \emph{Recall:} say the Gujarati for to applaud, then the Gujarati for pride, then say \emph{ṭukḍī} again
\end{itemize}

\begin{tcolorbox}[breakable,colback=teal!4,colframe=teal!35!black,title={Before you move on}]
What does \gu{ટુકડી} mean? (\textquotedblleft{}A squad; a small party (of police or soldiers)\textquotedblright{}.) Say it once more.
\end{tcolorbox}
\section[\texorpdfstring{\gu{જીવલેણ} (jīvleṇ)}{જીવલેણ (jīvleṇ)}]{\gu{જીવલેણ} (jīvleṇ) — deadly}

\label{lesson:GU-C310-jivlen}



\begin{quote}
Before the new one: say the Gujarati for pride, then the Gujarati for a squad.
\end{quote}

\subsection*{You'll want to know: \gu{જીવલેણ}}
\textbf{\gu{જીવલેણ}} — \emph{jīvleṇ} — \textquotedblleft{}deadly\textquotedblright{}.

\begin{quote}
\gu{આ} \gu{જંગલમાં} \gu{જીવલેણ} \gu{સાપ} \gu{છે}.
\end{quote}

\begin{quote}
\emph{ā jangalmā̃ jīvleṇ sāp che.}
\end{quote}

\begin{quote}
There are deadly snakes in this forest.
\end{quote}

\begin{cousinweb}
\textbf{\gu{જીવલેણ}} (\emph{jīvleṇ}) is built from \emph{jīv}, life, from Sanskrit \emph{jīvá}, \textquotedblleft{}living; life\textquotedblright{} (CDIAL 5239), and \emph{leṇ}, a taking, from \textbf{\gu{લેવું}}: what takes a life. \textbf{\gu{જીવલેણ} \gu{સાપ}}, a deadly snake.
\end{cousinweb}

\begin{grammarlens}[title={What you've built}]
One more word for telling what happened.
\end{grammarlens}

\subsection*{Your turn}
\begin{itemize}
\raggedright
  \item \emph{Say it:} \emph{jīvleṇ}
  \item \emph{Say it:} \emph{jīvleṇ}, once more
  \item \emph{Say it:} \emph{ā jangalmā̃ jīvleṇ sāp che.} — There are deadly snakes in this forest.
  \item \emph{Recall:} say the Gujarati for pride, then the Gujarati for a squad, then say \emph{jīvleṇ} again
\end{itemize}

\begin{tcolorbox}[breakable,colback=teal!4,colframe=teal!35!black,title={Before you move on}]
What does \gu{જીવલેણ} mean? (\textquotedblleft{}Deadly\textquotedblright{}.) Say it once more.
\end{tcolorbox}
\section[\texorpdfstring{\gu{સહાય} (sahāy)}{સહાય (sahāy)}]{\gu{સહાય} (sahāy) — assistance}

\label{lesson:GU-C310-sahay}



\begin{quote}
Before the new one: say the Gujarati for a squad, then the Gujarati for deadly.
\end{quote}

\subsection*{You'll want to know: \gu{સહાય}}
\textbf{\gu{સહાય}} — \emph{sahāy} — \textquotedblleft{}assistance\textquotedblright{}.

\begin{quote}
\gu{ટુકડીએ} \gu{ગામના} \gu{લોકોની} \gu{સહાય} \gu{લીધી}.
\end{quote}

\begin{quote}
\emph{ṭukḍīe gāmnā lokonī sahāy līdhī.}
\end{quote}

\begin{quote}
The squad took the villagers' help.
\end{quote}

A feminine noun.

\begin{cousinweb}
\textbf{\gu{સહાય}} (\emph{sahāy}) is a tatsama, Sanskrit \emph{sahāya}, \textquotedblleft{}a companion, a helper\textquotedblright{} (Monier-Williams); in Gujarati it is the help itself, a more formal word than \textbf{\gu{મદદ}}.
\end{cousinweb}

\begin{grammarlens}[title={What you've built}]
One more word for telling what happened.
\end{grammarlens}

\subsection*{Your turn}
\begin{itemize}
\raggedright
  \item \emph{Say it:} \emph{sahāy}
  \item \emph{Say it:} \emph{sahāy}, once more
  \item \emph{Say it:} \emph{ṭukḍīe gāmnā lokonī sahāy līdhī.} — The squad took the villagers' help.
  \item \emph{Recall:} say the Gujarati for a squad, then the Gujarati for deadly, then say \emph{sahāy} again
\end{itemize}

\begin{tcolorbox}[breakable,colback=teal!4,colframe=teal!35!black,title={Before you move on}]
What does \gu{સહાય} mean? (\textquotedblleft{}Assistance\textquotedblright{}.) Say it once more.
\end{tcolorbox}
\section[\texorpdfstring{\gu{મૂંઝાવું} (mū̃jhāvũ)}{મૂંઝાવું (mū̃jhāvũ)}]{\gu{મૂંઝાવું} (mū̃jhāvũ) — to be at a loss}

\label{lesson:GU-C310-mujhavu}



\begin{quote}
Before the new one: say the Gujarati for deadly, then the Gujarati for assistance.
\end{quote}

\subsection*{You'll want to know: \gu{મૂંઝાવું}}
\textbf{\gu{મૂંઝાવું}} — \emph{mū̃jhāvũ} — \textquotedblleft{}to be at a loss\textquotedblright{}.

\begin{quote}
\gu{જંગલમાં} \gu{બે} \gu{રસ્તા} \gu{જોઈને} \gu{ટુકડી} \gu{મૂંઝાઈ} \gu{ગઈ}.
\end{quote}

\begin{quote}
\emph{jangalmā̃ be rastā joīne ṭukḍī mū̃jhāī gaī.}
\end{quote}

\begin{quote}
When the squad saw two paths in the forest, it was at a loss.
\end{quote}

Past (he …): \gu{તે} \gu{મૂંઝાઈ} \gu{ગયો}, \emph{te mū̃jhāī gayo}.

\begin{cousinweb}
\textbf{\gu{મૂંઝાવું}} (\emph{mū̃jhāvũ}) is inherited from Sanskrit \emph{múhyati}, \textquotedblleft{}is bewildered\textquotedblright{}, and CDIAL lists Gujarati \emph{mujhāvũ}, \textquotedblleft{}to be perplexed\textquotedblright{} (CDIAL 10230). It is not knowing what to do next: \textbf{\gu{મૂંઝાઈ} \gu{ગઈ}}.
\end{cousinweb}

\begin{grammarlens}[title={What you've built}]
One more word for telling what happened.
\end{grammarlens}

\subsection*{Your turn}
\begin{itemize}
\raggedright
  \item \emph{Say it:} \emph{mū̃jhāvũ}
  \item \emph{Say it:} \emph{mū̃jhāvũ}, once more
  \item \emph{Say it:} \emph{jangalmā̃ be rastā joīne ṭukḍī mū̃jhāī gaī.} — When the squad saw two paths in the forest, it was at a loss.
  \item \emph{Recall:} say the Gujarati for deadly, then the Gujarati for assistance, then say \emph{mū̃jhāvũ} again
\end{itemize}

\begin{tcolorbox}[breakable,colback=teal!4,colframe=teal!35!black,title={Before you move on}]
What does \gu{મૂંઝાવું} mean? (\textquotedblleft{}To be at a loss\textquotedblright{}.) Say it once more.
\end{tcolorbox}
\section[\texorpdfstring{\gu{વિખેરાવું} (vikherāvũ)}{વિખેરાવું (vikherāvũ)}]{\gu{વિખેરાવું} (vikherāvũ) — to scatter; to disperse}

\label{lesson:GU-C310-vikheravu}



\begin{quote}
Before the new one: say the Gujarati for assistance, then the Gujarati for to be at a loss.
\end{quote}

\subsection*{You'll want to know: \gu{વિખેરાવું}}
\textbf{\gu{વિખેરાવું}} — \emph{vikherāvũ} — \textquotedblleft{}to scatter; to disperse\textquotedblright{}.

\begin{quote}
\gu{પોલીસ} \gu{આવી} \gu{ત્યારે} \gu{ટોળું} \gu{વિખેરાઈ} \gu{ગયું}.
\end{quote}

\begin{quote}
\emph{polīs āvī tyāre ṭoḷũ vikherāī gayũ.}
\end{quote}

\begin{quote}
When the police came, the crowd dispersed.
\end{quote}

Past (it …): \gu{તે} \gu{ટોળું} \gu{વિખેરાઈ} \gu{ગયું}, \emph{te ṭoḷũ vikherāī gayũ}.

\begin{cousinweb}
\textbf{\gu{વિખેરાવું}} (\emph{vikherāvũ}) is the intransitive partner of \emph{vikhervũ}, to scatter, which CDIAL lists under \emph{*viṣkirati}, a form it reconstructs (CDIAL 11985). A crowd or a cloud \textbf{\gu{વિખેરાઈ} \gu{જાય} \gu{છે}}: it breaks up.
\end{cousinweb}

\begin{grammarlens}[title={What you've built}]
One more word for telling what happened.
\end{grammarlens}

\subsection*{Your turn}
\begin{itemize}
\raggedright
  \item \emph{Say it:} \emph{vikherāvũ}
  \item \emph{Say it:} \emph{vikherāvũ}, once more
  \item \emph{Say it:} \emph{polīs āvī tyāre ṭoḷũ vikherāī gayũ.} — When the police came, the crowd dispersed.
  \item \emph{Recall:} say the Gujarati for assistance, then the Gujarati for to be at a loss, then say \emph{vikherāvũ} again
\end{itemize}

\begin{tcolorbox}[breakable,colback=teal!4,colframe=teal!35!black,title={Before you move on}]
What does \gu{વિખેરાવું} mean? (\textquotedblleft{}To scatter; to disperse\textquotedblright{}.) Say it once more.
\end{tcolorbox}
